\section{Introduction}
\label{sec:intro}

Every American is more directly governed by state law than by federal law.
State legislatures set property and income tax rates, fund public schools,
regulate police, determine criminal sentencing, and draw the very electoral
maps that shape future legislatures. Yet the algorithmic redistricting
literature has devoted nearly all of its energy to the 435 U.S.\ House seats.
State legislative chambers, which collectively contain roughly 7,400 seats
across 99 chambers in 50 states, have received far less systematic attention.

This imbalance is not accidental. Congressional redistricting is constitutionally
uniform: one person, one vote under \emph{Wesberry v.\ Sanders} (1964), a fixed
national total of 435, and a well-defined geography of 73,000 census tracts.
State legislative redistricting involves the same one-person-one-vote principle
(\emph{Reynolds v.\ Sims}, 1964), but the computational challenge is far greater:
district counts range from $k = 20$ (Alaska Senate) to $k = 400$ (New Hampshire
House), state-specific constitutional requirements impose varying contiguity and
compactness standards, and census tracts are often too coarse to support
meaningful redistricting at these district counts.

The F research track introduces a systematic program for applying the
\texttt{redist} pipeline---the minimum-edge-cut recursive bisection algorithm
of \citet{dellaLibera2026recursive}, implemented as a production Rust binary---to
all 50 state legislatures. This overview paper (F.0) frames the research program,
describes the computational infrastructure, previews the empirical findings,
and situates the F track within the broader portfolio.

\paragraph{Why state legislative redistricting matters more.}
Table~\ref{tab:chamber-scale} illustrates the stakes. Congressional gerrymandering
receives far more media attention, but state legislative gerrymandering determines
which party controls the body that enacts state policy, conducts criminal
trials, funds schools, and---critically---draws congressional maps.
In 2020, Republicans held trifectas (governorship plus both chambers) in 23 states;
Democrats in 15. These trifectas were themselves the product of state legislative
maps drawn after 2010 that are widely documented as among the most aggressive
partisan gerrymanders in American history
\citep{mcghee2014measuring, stephanopoulos2015partisan, chen2013unintentional}.
The F track asks: what would state legislative maps look like if drawn by the
same algorithmic principles that the B track applies to Congress?

\begin{table}[h]
\centering\small
\begin{tabular}{lrrrl}
\toprule
State & Chamber & Districts ($k$) & Tracts ($n$) & $k/n$ \\
\midrule
New Hampshire & House   & 400 & $\approx 320$ & 1.25 \\
California    & Assembly & 80  & $\approx 9{,}000$ & 0.009 \\
Texas         & House   & 150 & $\approx 5{,}250$ & 0.029 \\
Washington    & House   & 98  & $\approx 1{,}458$ & 0.067 \\
Wyoming       & House   & 60  & $\approx 140$ & 0.43 \\
Nebraska      & Legislature & 49 & $\approx 570$ & 0.086 \\
Alaska        & Senate  & 20  & $\approx 200$ & 0.10 \\
\bottomrule
\end{tabular}
\caption{Selected state chambers illustrating the range of district counts,
  tract counts, and the $k/n$ ratio that governs resolution selection.
  Values above the $0.05$ threshold (bold) require block-group resolution;
  New Hampshire and Wyoming require block-level redistricting.}
\label{tab:chamber-scale}
\end{table}

\paragraph{Organization.}
Section~\ref{sec:challenges} describes the two primary computational challenges:
the high-$k$ problem and the resolution problem. Section~\ref{sec:cli}
introduces the \texttt{redist} CLI extensions for state legislative redistricting.
Section~\ref{sec:overview} previews the six companion papers F.1--F.6.
Section~\ref{sec:findings} summarises the key empirical findings.
Section~\ref{sec:connections} situates the F track relative to the B and C tracks.
Section~\ref{sec:conclusion} concludes.
